% One-time imported draft. Editable native source; no Markdown regeneration.
\section{Make the delay slot explicit in the
  pre-test}\label{005-seven-segment-counter-p05-make-the-delay-slot-explicit-in-the-pre-test}

The zero-count branch must bypass the loop before any device
transaction. The retained \texttt{beq} instruction executes its
following delay-slot instruction for both outcomes. Placing \texttt{nop}
there makes the path easy to reason about: no lookup, write, or counter
change is hidden in the slot.

A post-test loop alone would execute its body once for an empty request.
If it then decremented zero, the count could become negative and repeat
far beyond the table. The initial test establishes that every entered
iteration has at least one requested output remaining.

The final \texttt{bne} uses the same explicit-slot rule. This chapter
keeps timing out of the output contract, but it still requires correct
architectural control flow. A delay slot is part of execution, not a
visual pause for a viewer.
